\scanpage{59}
\mspage{3}
% archive sheet number: 58
Leur conjonction est l'axiome des plus raisonnables
\[
W_\infty\subset W\subset W_0\ .
\]
Bien sûr, l'axiome L 5 a maintenant l'air un
peu étriqué, i.e. faiblard. On va voir très vite de
quoi on aura besoin ! \deleted{\unsure{Si}}

On veut donc définir un dérivateur
\[
\mathrm{HOT}_W
\]
à partir de la catégorie de modèles
\[
(\mathrm{Cat},W\subset\mathrm{Fl}(\mathrm{Cat})).
\]

Donc pour définir
\[
\mathrm{HOT}_W(S)\qquad (S\in\mathrm{Ob}\,\mathrm{Cat}),
\]
la recette générale est de prendre la catégorie
\[
M(S)=\mathrm{Cat}(S)=\underline{\mathrm{Hom}}(S^\circ,(\mathrm{Cat})),
\]
d'y introduire l'ens de flèches
\[
\left\{\begin{aligned}
&W_S\subset\mathrm{Fl}(M(S))\\
&u\in W_S\ \Longleftrightarrow\ u_s\in W\quad\forall s\in\mathrm{Ob}\,S\ .
\end{aligned}\right.
\]
On \unsure{doit|choisit} poser
\[
\mathrm{HOT}_W(S)\overset{\mathrm{def}}{=}\underline{\mathrm{Hom}}(S^\circ,\mathrm{Cat})W_S^{-1}.
\]
Mais la catégorie $\underline{\mathrm{Hom}}(S^\circ,\mathrm{Cat})$ s'interprète,
à \deleted{équivalence} \inserted{isomorphisme \unsure{can.|de c.|de cat.}} près, comme celle
\[
\underline{\mathrm{Fib.}}\underline{\mathrm{sc.}}(S)
\]
des catégories $X$ \deleted{fib.} fibrées sur $S$, munies
d'un ``scindage'' i.e. d'un choix des foncteurs
changement de base $u^*:X_t\longrightarrow X_s$ (pour $t\xrightarrow{u}s$
dans $S$), avec de plus \uline{transitivité} \uline{stricte}
\[
(vu)^*=u^*v^*\qquad\text{si}\quad t\mathrel{\overset{u}{-}}s\xrightarrow{v}r\text{ dans }S,
\]
